\begin{figure}[H]
\centering
\begin{tikzpicture}[font=\small,>=Stealth]

\filldraw[fill=blue!8,draw=blue,dashed,thick] (0,0)--(7,0)--(7,3)--(4.5,3)--(4.5,1.2)--(2.5,1.2)--(2.5,3)--(0,3)--cycle;
\draw[red!70!black,very thick] (1,2.3)--(1.5,.6)--(5.5,.6)--(6,2.3);
\fill (1,2.3) circle (2pt) node[above] {$p$};\fill (6,2.3) circle (2pt) node[above] {$q$};
\draw[teal,dashed] (1.5,.6) circle (.4);\draw[teal,dashed] (5.5,.6) circle (.4);
\node at (3.5,-.6) {small balls allow one more straight segment};

\end{tikzpicture}
\caption{A polygonal path can go around a gap in an open connected domain. The proof uses small convex balls to show that the reachable set and its complement are relatively open.}
\label{fig:polygonal-path}
\end{figure}
